\section{问题分析}
% 篇幅：**全节 ≤ 一页半，每个问题的分析 ≤ 3 行，每条创新点 3~5 行**。先一句总结性话语，再逐问分析，最后单列创新点。
%   标尺是参考件 `基于变物性耦合传热传质与动域模型的.pdf`：它的 2.1~2.4 各占 2~3 行、
%     创新点每条 4~5 行、全节 1.23 页。本模板排版下 1 行 ≈ 38 字、1 页 ≈ 989 字、一页半 ≈ 1480 字。
%   硬指标：每个「问题N的分析」≤ 3 行（≈114 字）；每条「创新点X：」**3~5 行**（≈114~190 字，
%     逐条量，三条各 5 行是正确写法；下限也是硬的 —— 只写一句时所有上限都满足，
%     于是"重写过"与"把上一版原样拿来用"分不出来）；全节 ≤ 1480 字（一页半），1250 字以上提醒。
%   **「一个数字都不写」只约束「问题N的分析」**：特征时间、下降百分比、倍数、量级、
%     参数明细全是结果，归第五节求解与第六节检验 —— 写在这里既超页、又和后面重复。
%     这里只回答三件事：这一问要什么、难在哪、走什么路线。
%   创新点可以引一个关键数字当证据（参考件每条都带一个）—— 它写的是
%     「做了什么 + 带来什么实质改进」，比问题分析长是对的。
%   `python lib/web/check_skeleton.py` 会逐条核（全节超一页半、或某节超行数上限都判 FAIL，
%     并列出分节明细告诉你该砍哪一段）。

% 下面「问题一/二/三的分析」是举例，不是固定三节：
%    题面有几个子问题，就写几节分析，一节不缺、一节不多。
%     例：题面四问就必须写到「问题四的分析」，漏掉任何一问都算漏节。
%     子问题数在 `5_problemN.tex` 的文件数上——建/删文件时同步建/删这里的分析节。
%     `python lib/web/check_skeleton.py` 会核这一条。

\subsection{问题一的分析}

\subsection{问题二的分析}

\subsection{问题三的分析}

\subsection{本文主要创新点}
% 逐条写，每条 **4~5 行**（≈190 字）：写清「做了什么 + 带来什么实质简化/改进」，
%   并引一个关键数字当证据（参考件每条都带一个）。不要写成方法罗列，也不要只写一句
%   —— 只写一句会短到 2 行，撑不起这一节。逐条独立成段，段间空行。
%   参考件是 3 条（每条 ≈178 字）；4 条时全节约 1215 字 ≈ 1.23 页，仍在一页半以内
%     —— **4~5 条也放得下**，不必为省页数把每条压成一句。

\textbf{创新点一：}

\textbf{创新点二：}

\textbf{创新点三：}
